% Terminal apparatus of the concise study: the scope note and the References
% for the sources this companion actually uses, at the loci it cites. The
% revision timestamp and the rights colophon follow from the shared
% generation record.

\runninghead{Scope and References}
\section{Appendix: Scope and Qualifications}

{\footnotesize
\textbf{Edition, Mass and occurrence.} The Mass treated is \latin{Dominica
quinta quae superfuit post Epiphaniam}, II classis, of the \work{Missale
Romanum}, editio typica, Typis Polyglottis Vaticanis 1962, printed
pp.~420--421, nos.~1749--1758, as said on Sunday 8~November 2026, the
twenty-fourth Sunday after Pentecost, in the universal 1962 calendar. Its place in 2026, the omission of the commemoration of
the Four Crowned Martyrs, the colour, the Gloria, the Credo and the Trinity
Preface rest on the Missal's general rubrics and on its printed direction
after the Secret, as \texttt{research/context.md} records. The Ordinary and
the Office are not treated.

\textbf{Text and English.} The Latin of the ten texts was read on page images
of the CMAA facsimile of the 1962 typical edition and collated with the
Ratisbon printing of 1862 (Pustet); the collation is in
\texttt{propers/verified.md}; the expansive study prints every text whole.
The English of the chants and orations is that of the Lasance \work{New Roman
Missal}
(Benziger, 1945), public domain in the United States by non-renewal; the
lessons are quoted in the Douay--Rheims in Challoner's revision. None of this
English is an approved liturgical translation of the 1962 Missal. The Fathers are quoted in the English
translations named below or in their Latin with an English rendering made for
the expansive study; no English is composed for a liturgical text.

\textbf{Reception.} The three interpretations are developed in full in the
expansive study; their loci, editions and bounds are recorded in
\texttt{research/scope.md} and \texttt{research/interpretations.md}. Every
connection drawn between two texts of the Mass belongs to the interpretations
unless an author is named as making it. The Latin of Jerome, Hilary, Rabanus
and Ambrose was read on page images; Aquinas, Bede, Cassiodorus, Honorius, Rupert,
Durandus and Sicard in uncorrected optical or web text; Chrysostom on the
Psalms, Theodoret and Theophylact in uncorrected Greek or Latin optical text,
so that Chrysostom on Psalms~43 and 129 is given in paraphrase. Augustine's work against
Parmenian is cited only as Aquinas cites it. Honorius, Cassiodorus, Theodoret, Theophylact, Rupert, Durandus and Sicard are
ecclesiastical writers and are cited in support only. St Thomas Aquinas's
conclusion on the punishment of heretics is reported as his; the later
development of the Church's teaching on religious liberty is not treated.

\textbf{Chronology.} Every biblical date on page~2 is generated from the
repository's Scripture chronology corpus; the expansive study's Scriptural
Date and Location appendix gives each passage's sources and evidence. The
record holds no date for the occasion of Psalms~43, 84
or 129 or for the telling of the parable.

\textbf{Rights and use.} The Latin is published on 17 U.S.C. 103(b) together
with its public-domain witness, the Pustet printing of 1862. The Douay--Rheims
and the Lasance English are quoted from registered editions. The other
quotations are short and come from editions published before 1931, named
below, or from medieval and ancient Latin; the Corpus Thomisticum and
Wikisource texts are used only for Latin that is itself medieval or
ancient. Scripture, liturgical texts and quoted translations are
not relicensed by the repository's licence. This is a study aid to the 1962
Missal, not an official liturgical book, translation or critical edition.
\par}

\section{References}

{\footnotesize
\begin{itemize}[leftmargin=1.2em,itemsep=.15em]
\item \work{Missale Romanum ex decreto Sacrosancti Concilii Tridentini
restitutum, Summorum Pontificum cura recognitum}, editio typica (Typis
Polyglottis Vaticanis, 1962), CMAA facsimile: \work{Rubricae generales} and
\work{Rubricae generales Missalis romani}; Proprium de Tempore,
pp.~420--421, nos.~1749--1758.
\item \work{Missale Romanum} (Ratisbon: Pustet, 1862), pp.~42--44 and
357--358.
\item Lasance, F.~X., and Walsh, F.~A., \work{The New Roman Missal in Latin and
English}, revised edition (New York: Benziger Brothers, 1945), pp.~196 and
744--747.
\item \work{The Holy Bible Translated from the Latin Vulgate} (Douay--Rheims,
Challoner revision), Project Gutenberg text: Gen 50:17, 19;
Ps 43:7--10, 22; 84:2; 129:1--4; Jer 29:1--14; Mt 7:7; 13:1--3, 24--39;
Mk 11:20--25; Rom 8:36; 1~Cor 5:13; Col 3:1--17; 2~Pet 2:8.
\item Ambrose, \work{Enarrationes in XII Psalmos}, on Ps 43, nn.~21--22, PL~14
(Paris, 1845), col.~1154.
\item Augustine, \work{Sermo} 73 (Sermon XXIII in NPNF), first series,
vol.~6; \work{Epistula} 93 to Vincentius, NPNF, first series, vol.~1;
\work{Enarrationes in Psalmos} 42, 43, 84 and 129, NPNF, first series,
vol.~8.
\item Bede, \work{In Marci evangelium expositio} III, on 11:22--25, Migne
Latin text (Wikisource).
\item Bellarmine, Robert, \work{Explanatio in Psalmos}, in John O'Sullivan's
English, \work{A Commentary on the Book of Psalms}, on Pss 43, 84 and 129.
\item Cassiodorus, \work{Expositio Psalmorum}, on Ps 43:8, PL~70.
\item Chrysostom, John, \work{Homilies on Matthew} 46, NPNF, first series,
vol.~10; \work{Homilies on Colossians} 8 and 9, NPNF, first series, vol.~13;
\work{Expositio in Psalmos}, on Pss 43 and 129, PG~55.
\item Hilary of Poitiers, \work{Tractatus super Psalmos}, on Ps 129,
nn.~2--3, PL~9 (Paris, 1844), cols.~719--720.
\item Honorius Augustodunensis, \work{Speculum Ecclesiae}, the sermon on the
cockle headed \latin{Dominica tercia}, PL~172.
\item Jerome, \work{Commentarii in Matthaeum} II, on 13:36--43, PL~26 (Paris,
1845), cols.~93--94; \work{In Hieremiam} V, on 29:5--14, PL~24 (Paris, 1865),
cols.~891--893.
\item Rabanus Maurus, \work{Commentaria in Matthaeum} IV, on 13:24--30, PL~107
(Paris, 1864), cols.~947--948.
\item Theodoret of Cyrus, \work{Interpretatio epistolae ad Colossenses}, on
3:15, PG~82; Theophylact of Ohrid, \work{Expositio in epistolam ad
Colossenses}, on 3:15, PG~124, Latin version, cols.~1261--1264.
\item Thomas Aquinas, \work{Super Evangelium S.~Matthaei lectura}, cap.~13, in
\work{Opera omnia}, vol.~3 (Venice: Bettinelli, 1745), pp.~183--185;
\work{Super Psalmos}, on Ps 43, n.~3, Corpus Thomisticum text and
\work{Opera omnia}, vol.~14 (Parma, 1863); \work{Super epistolam ad
Colossenses}, cap.~3, lect.~3, Corpus Thomisticum text and \work{In omnes
D.~Pauli Apostoli epistolas commentaria}, vol.~2 (Liège: Dessain, 1857);
\work{Summa theologiae} II-II q.~10 a.~8 and q.~11 a.~3, in the English
Dominican translation.
\item Durandus, William, \work{Rationale divinorum officiorum} VI.140, in the
Lyon printing of 1612, tomus secundus, p.~412.
\item Rupert of Deutz, \work{De divinis officiis} XII.23, PL~170,
cols.~329--330.
\item Sicard of Cremona, \work{Mitrale} VI.23, PL~213.
\end{itemize}
\par}
